\section{Evaluation Protocol and Comparisons}
\label{app:experimental}
\subsection{Tasks, Interface, and Execution Settings}
\label{app:configuration}
The training design partitions 3,553 gamefiles into 2,843 TRAIN\_UPDATE,
355 TRAIN\_SELECT, and 355 TRAIN\_AUDIT tasks. The first pool supplies
rollouts, analysis, repair experiments, and training; the second supplies
candidate-selection results; the third supports research-role qualification.
All states, seeds, and branches belonging to one gamefile inherit its pool.
Final seen and unseen evaluations remain outside RePair's adaptive research
and selection loop.


\begin{table}[!htbp]
\caption{Executed dual-view validation and evaluation settings. Repair
verification and policy selection use different seed protocols.}
\label{tab:config}\centering\small
\begin{tabularx}{\linewidth}{L{1.6in}X}\toprule
Field & Setting\\\midrule
Base model & Qwen2.5-3B-Instruct; retained parent adapter continued, base weights frozen.\\
Research models & External GPT API; recent validation records \texttt{gpt-5.6-sol}; no teacher takeover during branch execution.\\
Training-side pools & UPDATE 2,843; SELECT 355; AUDIT 355, disjoint by gamefile.\\
Final benchmark & Unseen 134 and seen 140, kept separate from selection.\\
Interface and budget & Strict single-action JSON; full menu; eight executed transitions; 30 environment steps; 60 policy attempts; three consecutive nonexecutions stop.\\
Selection / repair seeds & Selection: 17; repair: 17, 31, 47, 73, 101.\\
LoRA & Rank 16; alpha 32; dropout 0.05; no bias; attention projections and gate/up/down MLP projections.\\
Optimizer & AdamW; learning rate $2\!\times\!10^{-5}$; betas $(0.9,0.999)$; epsilon $10^{-8}$; weight decay 0.01; gradient norm cap 1.0.\\
Schedule and batches & Linear schedule; one warmup step; microbatch 1; accumulation 4; four passes; eight optimizer steps.\\
Data and randomness & Eight examples; max length 1,248 without truncation; training and data seed 17; one trained run.\\
Loss and precision & Assistant-only token mean per accumulation group; BF16; one GPU; no packing or distributed reduction.\\
Checkpoint & Final update only; no intermediate evaluation or checkpoint selection.\\
Statistics & Exact task-paired counts; no significance or multi-training-seed estimate.\\\bottomrule
\end{tabularx}
\end{table}

\paragraph{Shared raw prompt.}
The seven-line prompt begins with a fixed format identifier, followed by
the task goal, current observation, executed-transition history, complete
admissible-action menu, interface feedback, and output requirement. The
context fields are JSON-serialized, and the menu retains its original order.
The output requirement is \texttt{\{"action":"<command>"\}}.
The parser accepts exactly one top-level JSON object with one string member
named \texttt{action}. It removes outer transport whitespace and ASCII
spaces at action edges, then tests exact case-sensitive menu membership.
Duplicate members, trailing data, nonstandard constants, multiline or control
characters, and actions longer than 256 codepoints fail parsing.
No action-repair or canonicalization layer is applied.

\paragraph{Budget accounting.}
A format failure or parsed out-of-menu action consumes one policy attempt
and zero environment steps. An admissible action consumes one of each;
successful environment execution resets the consecutive-nonexecution count.
After a nonexecution, the consecutive-three limit takes precedence over
policy-attempt exhaustion. After execution, termination priority is
infrastructure failure, environment terminal status, the 30-step limit,
then the 60-attempt limit. Each episode retains its termination reason.

\paragraph{Frozen evaluation assets.}
The shared evaluation protocol and the two task lists are fixed before
evaluation and kept unchanged across methods. The lists contain 140 seen
games and 134 unseen games, respectively. Version records and file checksums
identify these assets, and execution records link each run to the assets it
used. Full task lists and run records are retained with the experiment
artifacts; the accompanying numerical supplement preserves the table inputs and
result provenance without private server paths.
The shared external interface and action budget coexist with model-specific
inference profiles. In particular, high-reasoning API inference and local
128-token generation have different computation costs. The shared action budget is not a matched inference-compute budget;
we do not claim equal token cost or latency across these methods.

\subsection{Reference Checkpoints on the Final Benchmark Splits}
\label{app:aggregation}
The benchmark columns are Pick, Clean, Cool, Look, Heat, Pick2, and All.
They correspond to pick-and-place, clean-and-place, cool-and-place,
look-at-object-in-light, heat-and-place, and pick-two-object tasks.
Table~\ref{tab:main} reports unseen-134; Table~\ref{tab:seen} uses the same
method roster and columns for seen-140. These reference measurements give
task context across heterogeneous model scales and inference settings.
The RePair candidate has not been evaluated on either final split;
the tables therefore contain reference models only.
\begin{table}[H]
\caption{\textbf{Unseen ALFWorld reference results.} Success (\%) on 134 games using seed 17. All is task-weighted. These seven reference models use the shared task and environment protocol with model-specific inference settings. They provide task context, not a matched comparison with RePair; its final benchmark evaluation remains incomplete.}
\label{tab:main}
\centering
\small
\setlength{\tabcolsep}{4pt}
\begin{tabularx}{\linewidth}{@{}Xrrrrrrr@{}}
\toprule
Method & Pick & Clean & Cool & Look & Heat & Pick2 & All\\
\midrule
GPT 5.6 sol high & \res{bench-unseen-strong-pick} & \res{bench-unseen-strong-clean} & \res{bench-unseen-strong-cool} & \res{bench-unseen-strong-look} & \res{bench-unseen-strong-heat} & \res{bench-unseen-strong-pick2} & \res{bench-unseen-strong-all}\\
\midrule
Qwen2.5-3B-Instruct (base model) & \res{bench-unseen-base-pick} & \res{bench-unseen-base-clean} & \res{bench-unseen-base-cool} & \res{bench-unseen-base-look} & \res{bench-unseen-base-heat} & \res{bench-unseen-base-pick2} & \res{bench-unseen-base-all}\\
\midrule
SkillRL & \res{bench-unseen-skillrl-pick} & \res{bench-unseen-skillrl-clean} & \res{bench-unseen-skillrl-cool} & \res{bench-unseen-skillrl-look} & \res{bench-unseen-skillrl-heat} & \res{bench-unseen-skillrl-pick2} & \res{bench-unseen-skillrl-all}\\
GiGPO & \res{bench-unseen-gigpo-pick} & \res{bench-unseen-gigpo-clean} & \res{bench-unseen-gigpo-cool} & \res{bench-unseen-gigpo-look} & \res{bench-unseen-gigpo-heat} & \res{bench-unseen-gigpo-pick2} & \res{bench-unseen-gigpo-all}\\
SEED & \res{bench-unseen-seed-pick} & \res{bench-unseen-seed-clean} & \res{bench-unseen-seed-cool} & \res{bench-unseen-seed-look} & \res{bench-unseen-seed-heat} & \res{bench-unseen-seed-pick2} & \res{bench-unseen-seed-all}\\
Agent\_G2 & \res{bench-unseen-agentg2-pick} & \res{bench-unseen-agentg2-clean} & \res{bench-unseen-agentg2-cool} & \res{bench-unseen-agentg2-look} & \res{bench-unseen-agentg2-heat} & \res{bench-unseen-agentg2-pick2} & \res{bench-unseen-agentg2-all}\\
OPID & \res{bench-unseen-opid-pick} & \res{bench-unseen-opid-clean} & \res{bench-unseen-opid-cool} & \res{bench-unseen-opid-look} & \res{bench-unseen-opid-heat} & \res{bench-unseen-opid-pick2} & \res{bench-unseen-opid-all}\\
\bottomrule
\end{tabularx}
\end{table}

\begin{table}[H]
\caption{\textbf{Seen ALFWorld reference results.} Success (\%) on 140 games using seed 17. All is task-weighted. These seven reference models use the shared task and environment protocol with model-specific inference settings. They provide task context, not a matched comparison with RePair; its final benchmark evaluation remains incomplete.}
\label{tab:seen}
\centering
\small
\setlength{\tabcolsep}{4pt}
\begin{tabularx}{\linewidth}{@{}Xrrrrrrr@{}}
\toprule
Method & Pick & Clean & Cool & Look & Heat & Pick2 & All\\
\midrule
GPT 5.6 sol high & \res{bench-seen-strong-pick} & \res{bench-seen-strong-clean} & \res{bench-seen-strong-cool} & \res{bench-seen-strong-look} & \res{bench-seen-strong-heat} & \res{bench-seen-strong-pick2} & \res{bench-seen-strong-all}\\
\midrule
Qwen2.5-3B-Instruct (base model) & \res{bench-seen-base-pick} & \res{bench-seen-base-clean} & \res{bench-seen-base-cool} & \res{bench-seen-base-look} & \res{bench-seen-base-heat} & \res{bench-seen-base-pick2} & \res{bench-seen-base-all}\\
\midrule
SkillRL & \res{bench-seen-skillrl-pick} & \res{bench-seen-skillrl-clean} & \res{bench-seen-skillrl-cool} & \res{bench-seen-skillrl-look} & \res{bench-seen-skillrl-heat} & \res{bench-seen-skillrl-pick2} & \res{bench-seen-skillrl-all}\\
GiGPO & \res{bench-seen-gigpo-pick} & \res{bench-seen-gigpo-clean} & \res{bench-seen-gigpo-cool} & \res{bench-seen-gigpo-look} & \res{bench-seen-gigpo-heat} & \res{bench-seen-gigpo-pick2} & \res{bench-seen-gigpo-all}\\
SEED & \res{bench-seen-seed-pick} & \res{bench-seen-seed-clean} & \res{bench-seen-seed-cool} & \res{bench-seen-seed-look} & \res{bench-seen-seed-heat} & \res{bench-seen-seed-pick2} & \res{bench-seen-seed-all}\\
Agent\_G2 & \res{bench-seen-agentg2-pick} & \res{bench-seen-agentg2-clean} & \res{bench-seen-agentg2-cool} & \res{bench-seen-agentg2-look} & \res{bench-seen-agentg2-heat} & \res{bench-seen-agentg2-pick2} & \res{bench-seen-agentg2-all}\\
OPID & \res{bench-seen-opid-pick} & \res{bench-seen-opid-clean} & \res{bench-seen-opid-cool} & \res{bench-seen-opid-look} & \res{bench-seen-opid-heat} & \res{bench-seen-opid-pick2} & \res{bench-seen-opid-all}\\
\bottomrule
\end{tabularx}
\end{table}


For disjoint task-family subsets $\mathcal D_c$ within one split,
\begin{equation}
\widehat{\SR}_c=\frac{1}{N_c}\sum_{t\in\mathcal D_c}\bar Y_t,\qquad
\widehat{\SR}_{\mathrm{All}}=\frac{\sum_c N_c\widehat{\SR}_c}{\sum_c N_c},\qquad
\bar Y_t=\frac{1}{J_{\mathrm{eval}}}\sum_jY_{tj},\quad N_c=|\mathcal D_c|.
\label{eq:task-types}
\end{equation}
Thus All weights families by their task counts. Tables report
$100\widehat{\SR}$ as percentages. The primary benchmark uses
$J_{\mathrm{eval}}=1$. Raw successes, family counts, and their mapping to
ALFWorld task IDs are exported with each table. Seen and unseen have separate
denominators throughout.


\subsection{Scope of Internal Comparisons}
\label{app:comparisons}
The framework supports simple versus hierarchical analysis, historical
memory access, alternative Planner rankings, and action-only versus
dual-view supervision. These define useful controlled experiments, but
the present development sequence changes multiple components and does
not instantiate those matched comparisons. In particular, there is no
completed weaker-filter training control or equal-compute action-only
control for the reported candidate. We do not report a component effect
from the difference between sequential portfolios.

\subsection{Repair Yield, Cost, and Statistical Units}
\label{app:metrics}
A registered-unit discovery-yield comparison would require a common
failure universe, proposal budget, and verifier. The current report instead
uses the directly observed number of tested states and branches in each
portfolio. These denominators exclude unselected candidates and therefore
must not be interpreted as a matched discovery-yield comparison.
Cost per useful unit is
\begin{equation}
C_{\mathrm{repair}}(a)
=\frac{N_{\mathrm{branches}}(a)}{N_{\mathrm{Benefit\ units}}(a)}.
\label{eq:cost}
\end{equation}
A zero-Benefit run reports its total cost and an undefined ratio.
The reported branch counts include complete unsuccessful tests. They do
not constitute total campaign cost: upstream API calls, retries, engineering
recovery, and training compute have not been reconciled into a comparable
end-to-end cost estimate.

Among candidates with complete formal verdicts, the Benefit share is
\begin{equation}
p_B=\frac{N_B}{N_{\mathrm{tested}}},\qquad
N_{\mathrm{tested}}=N_B+N_H+N_{NS}+N_{NF}+N_U.
\label{eq:precision}
\end{equation}
This candidate-level denominator differs from a discovery yield measured
over a common registered failure universe. The outcome breakdown retains neutral and uncertain cases alongside
Harm, while unbound, unsupported, invalid, and incomplete proposals have
separate counts.


The 355 paired selection tasks are the statistical units for the final
validation. The paired outcome table has 2 shared successes, 1
candidate-only success, 1 parent-only success, and 351 shared failures.
We provide exact counts rather than a significance claim based on these
two discordant tasks. Five continuation seeds at a repair source are
repetitions nested within that source, not five independent repairs.

\subsection{The Verified-Action Control}
\label{app:cso-control}
A matched action-only control would share the parent, source evidence,
action coverage, and evaluation protocol while omitting auxiliary targets.
CSO uses verified alternatives with a different critical-step procedure
and DPO objective \citep{li2026cso}; such an internal RePair control would
not constitute a CSO reproduction. This control has not been run here.

\subsection{Released Checkpoints in the Reference Benchmark}
\label{app:baseline-status}
The external comparison uses the five author-released checkpoints in
Table~\ref{tab:checkpoint-identities}. Their selection precedes shared
benchmark evaluation. Each is loaded at the previously frozen repository
revision; the experimental release records that revision, conversion steps,
weight hashes, tokenizer, and episode outputs. The model names and nominal
sizes identify a heterogeneous checkpoint comparison.

\begin{table}[htbp]
\caption{Author-released ALFWorld checkpoints evaluated under the common
RAW\_WITH\_MENU interface. Repository names identify the selected releases; immutable revisions,
conversion receipts, and completion records are retained with the experimental
artifacts.}
\label{tab:checkpoint-identities}
\centering\small
\begin{tabularx}{\linewidth}{@{}L{0.80in}L{0.42in}X@{}}
\toprule
Method & Scale & Checkpoint repository\\\midrule
SkillRL & 7B & \nolinkurl{Jianwen/Alfworld-7B-RL}\\
GiGPO & 7B & \nolinkurl{langfeng01/GiGPO-Qwen2.5-7B-Instruct-ALFWorld}\\
SEED & 3B & \nolinkurl{Jinyang23/Seed-AlfWorld-3B}\\
Agent\_G2 & 7B & \nolinkurl{xiamoent/Agent-G2-alfworld-7b}\\
OPID & 1.7B & \nolinkurl{Jinyang23/OPID-ALFWorld-1.7B}\\
\bottomrule
\end{tabularx}
\end{table}

All five use the shared ALFWorld task/environment contract and primary seed
slot, while method-faithful adapters preserve inference behavior required by
the released checkpoint. Project RePair Memory and Harness are disabled. A
frozen native component is retained when it is part of the released method's
inference definition---for example, SkillRL uses its frozen official SkillBank
without benchmark-time writeback or dynamic updating. The evaluation performs
no checkpoint adaptation from benchmark feedback. Model scales, checkpoint
identities, decode profiles, and adapter choices are recorded with the run.
The reported parent--candidate selection comparison uses one lineage,
but it does not isolate individual algorithmic components.

SkillRL's sharded actor checkpoint requires its recorded conversion to a
loadable model; no optimizer state is used for benchmark-time learning.
GiGPO, SEED, Agent\_G2, and OPID likewise retain the chosen release identity
through loading and evaluation. The final export accounts for all 274
primary cells per checkpoint, with 134 unseen and 140 seen tasks, and
records raw responses, parser outcomes, actions, steps, termination, tokens,
and latency. Aggregate scores are derived after completion. The completed task-level exports for all five released baselines are the
source of Tables~\ref{tab:main} and~\ref{tab:seen}.
